\exercisesection{4}

\begin{enumerate}
\def\labelenumi{\arabic{enumi}.}
\tightlist
\item
  Let A be a commutative ring which is Hausdorff and complete with a filtration \((a_{n})_{n\geq0}\) such that \(a_{0}=A\) . Let M, \(M'\) , N be three A-modules, each with the filtration induced by that on A and the topology defined by that filtration; we write \(\overline{M}=M/a_{1}M\) , \(\overline{M}'=M'/a_{1}M'\) , \(\overline{N}=N/a_{1}N\) . Let \(f: M\times M'\to N\) be a bilinear mapping and \(\bar{f}: \overline{M}\times\overline{M'}\to\overline{N}\) the \((A/a_{1})\) -bilinear mapping derived from f by taking quotients. Let \(y\in N\) , \(\bar{x}\in\overline{M}\) , \(\bar{x}'\in\overline{M'}\) be such that: (1) \(\bar{f}(\bar{x},\bar{x}')\) is the class \(\bar{y}\) of y in \(\overline{N}\) ; (2) every element of \(\overline{N}\) can be written as
\end{enumerate}

\[
\bar {f} (\overline {{{x}}}, \overline {{{z}}} ^ {\prime}) + \bar {f} (\overline {{{z}}}, \overline {{{x}}} ^ {\prime}),
\]

where \(\bar{z} \in \overline{\mathbf{M}}\) and \(\bar{z}' \in \overline{\mathbf{M}}'\) . Show that, if N is Hausdorff and M and M' are complete, there exists \(x \in \bar{x}\) and \(x' \in \bar{x}'\) such that \(f(x, x') = y\) (argue by induction as in the proof of Hensel's lemma). Under what conditions are \(x\) and \(x'\) determined uniquely?

\begin{enumerate}
\def\labelenumi{\arabic{enumi}.}
\setcounter{enumi}{1}
\tightlist
\item
  Let A be a local ring, m its maximal ideal, k = A/m its residue field and \(f: A \to k\) the canonical homomorphism. Let \(P \in A[X]\) be a monic polynomial of degree n.~We write \(B = A[X]/P \cdot A[X]\) and denote by x the canonical image of X in B.
\end{enumerate}

\begin{enumerate}
\def\labelenumi{(\alph{enumi})}
\item
  Let \(\mathbf{Q}, \mathbf{Q}'\) be two strongly relatively prime monic polynomials in \(\mathbf{A}[\mathbf{X}]\) such that \(\mathbf{P} = \mathbf{QQ}'\) . Show that the ring \(\mathbf{B}\) is the direct sum of the ideals \(\mathbf{B}. \mathbf{Q}(x)\) and \(\mathbf{B}. \mathbf{Q}'(x)\) .
\item
  Conversely, let \(B = b \oplus b'\) be a decomposition of B as a direct sum of two ideals. Show that there exist polynomials Q, \(Q'\) in A{[}X{]} satisfying the hypotheses of (a) and such that \(\mathfrak{b} = \mathrm{B}.Q(x)\) , \(\mathfrak{b}' = \mathrm{B}.Q'(x)\) . (Show first that b/mb and \(b'/mb'\) are generated by the images in B/mB of monic polynomials \(Q_{0} \in \bar{f}^{-1}(\mathfrak{b})\) , \(Q'_{0} \in \bar{f}^{-1}(\mathfrak{b}')\) such that
\end{enumerate}

\[
\bar {f} (\mathrm{P}) = \bar {f} (\mathrm{Q} _ {0}) \bar {f} (\mathrm{Q} _ {0} ^ {\prime}).
\]

Let \(r = \deg(Q_0)\) , \(s = \deg(Q_0')\) . Show that b is a free A-module with basis \(Q_0(x)\) , \(xQ_0(x)\) , \(\ldots\) , \(x^{s-1}Q_0(x)\) (apply Chapter II, § 3, no. 2, Proposition 5); we may then write \(x^sQ_0(x) = a_0Q_0(x) + a_1xQ_0(x) + \cdots + a_{s-1}x^{s-1}Q_0(x)\) , where \(a_i \in A\) ( \(0 \leqslant i \leqslant s-1\) ); show that, if we write

\[
\mathrm{Q} ^ {\prime} (\mathrm{X}) = \mathrm{X} ^ {s} - (a _ {0} + a _ {1} \mathrm{X} + \dots + a _ {s - 1} \mathrm{X} ^ {s - 1}),
\]

then \(\bar{f}(\mathbf{P}) = \bar{f}(\mathbf{Q}_0)\bar{f}(\mathbf{Q}')\) and \(\mathbf{Q}'\in \bar{f}^{-1}(\mathfrak{b}')\) . Similarly define \(\mathbf{Q}\) starting with \(\mathbf{Q}'\) and \(\mathbf{Q}_0\) and show that \(\mathbf{Q}\) and \(\mathbf{Q}'\) solve the problem, using Proposition 2 of no. 1.)

¶ 3. Let A be a local ring, m its maximal ideal, k = A/m its residue field and f: A → k the canonical homomorphism. Show that the two following conditions are equivalent:

\begin{enumerate}
\def\labelenumi{(\Alph{enumi})}
\setcounter{enumi}{7}
\item[(H)]
  For every monic polynomial \(\mathbf{P} \in \mathbf{A}[\mathbf{X}]\) and every decomposition of \(\bar{f}(\mathbf{P}) \in k[\mathbf{X}]\) as a product \(\bar{f}(\mathbf{P}) = \overline{\mathbf{Q}}\) . \(\overline{\mathbf{Q}}'\) of relatively prime monic polynomials, there exist two monic polynomials \(\mathbf{Q}\) , \(\mathbf{Q}'\) in \(\mathbf{A}[\mathbf{X}]\) such that \(\bar{f}(\mathbf{Q}) = \overline{\mathbf{Q}}\) , \(\bar{f}(\mathbf{Q}') = \overline{\mathbf{Q}}'\) and \(\mathbf{P} = \mathbf{QQ'}\) .
\item[(C)]
  Every commutative algebra over A which is a finitely generated A-module is a direct composition of A-algebras which are local rings.
\end{enumerate}

(To prove that (H) implies (C), argue as in Proposition 8 of no. 6, using Exercise 2 (a). To see that (C) implies (H), show first that, for every commutative A-algebra B which is a finitely generated A-module, every decomposition of B/mB as a direct sum of two ideals is necessarily of the form b/mb ⊕ b'/mb', where B = b ⊕ b' is a decomposition of B as a direct sum of two ideals; then use Exercise 2 (b).)

A local ring satisfying conditions (H) and (C) is called Henselian. Every complete Hausdorff local ring is Henselian. If A is Henselian and B is a commutative A-algebra which is a local ring and a finitely generated A-module, then B is Henselian.

\begin{enumerate}
\def\labelenumi{\arabic{enumi}.}
\setcounter{enumi}{3}
\tightlist
\item
  \begin{enumerate}
  \def\labelenumii{(\alph{enumii})}
  \tightlist
  \item
    Let \((\mathbf{A}_{\alpha}, \phi_{\alpha \beta})\) be a direct system of Henselian local rings, the homomorphisms \(\phi_{\alpha \beta}\) being local. Show that the local ring \(\mathbf{A} = \varinjlim \mathbf{A}_{\alpha}\) (Chapter II, § 3, Exercise 16) is Henselian (use criterion (H) of Exercise 3).
  \end{enumerate}
\end{enumerate}

\begin{enumerate}
\def\labelenumi{(\alph{enumi})}
\setcounter{enumi}{1}
\tightlist
\item
  Let K be a commutative field and L an algebraic extension of K. Deduce from (a) that the ring \(L \otimes_{K} K[[X_{1}, \ldots, X_{n}]]\) is Henselian. * (c) Let A be a Henselian local ring and B a commutative A-algebra which is integral over A (Chapter V) and is a local ring. Show that B is a Henselian ring (use (a)). *
\end{enumerate}

¶ 5. Let A be a Henselian local ring, B a (not necessarily commutative) A-algebra which is finitely generated A-module, b a two-sided ideal of B and let \(\overline{\mathbf{B}} = \mathbf{B} / \mathfrak{b}\) .

\begin{enumerate}
\def\labelenumi{(\alph{enumi})}
\item
  Show that every idempotent \(\varepsilon\) in \(\overline{\mathbf{B}}\) is the canonical image of an idempotent in \(\mathbf{B}\) (reduce it to the commutative case, considering the subalgebra of \(\mathbf{B}\) generated by a single element).
\item
  Let \((\varepsilon_{n})_{n\geqslant 1}\) be an infinite sequence of elements of \(\overline{\mathbf{B}}\) such that
\end{enumerate}

\[
\varepsilon_ {i} \varepsilon_ {j} = \delta_ {i j} \varepsilon_ {j}
\]

for every ordered pair of indices \((i,j)\) . Show that there exists in B an orthogonal sequence \((e_n)_{n\geq 1}\) of idempotents such that \(\varepsilon_{n}\) is the canonical image of \(e_n\) for all \(n\) . (Argue by induction as in Algebra, Chapter VIII, §6, Exercise 10; observe that, if \(e,e^{\prime}\) are two idempotents of B such that \(ee^{\prime} = 0\) , \(e^{\prime} - e^{\prime}e = e''\) is an idempotent such that \(ee'' = e''e = 0\) .)

\begin{enumerate}
\def\labelenumi{(\alph{enumi})}
\setcounter{enumi}{2}
\tightlist
\item
  Suppose now that b is the Jacobson radical of B. Let n be an integer and \((\varepsilon_{ij})\) ( \(1 \leqslant i \leqslant n, 1 \leqslant j \leqslant n\) ) a family of elements of \(\overline{B}\) such that \(\varepsilon_{ij}\varepsilon_{hk} = \delta_{jh}\varepsilon_{ik}\) and \(1 = \sum_{i=1}^{n} \varepsilon_{ii}\) . Show that there exists in B a family \((e_{ij})\) ( \(1 \leqslant i \leqslant n, 1 \leqslant j \leqslant n\) ) such that \(e_{ij}e_{hk} = \delta_{jh}e_{ik}\) and \(1 = \sum_{i=1}^{n} e_{ii}\) and that \(\varepsilon_{ij}\) is the canonical image of \(e_{ij}\) for every ordered pair of indices. (Use (b), Exercise 11 of Algebra, Chapter VIII, §6, and Exercise 9 of Algebra, Chapter VIII, §1.) Deduce that, if \(\overline{B}\) is isomorphic to a matrix ring \(\mathbf{M}_{r}(\overline{\mathbf{D}})\) , where \(\overline{D}\) is a not necessarily commutative field, then B is isomorphic to a matrix ring \(\mathbf{M}_{r}(\mathbf{D})\) , where D is an A-algebra which is a finitely generated A-module and whose Jacobson radical \(\mathfrak{d}\) is such that \(D/\mathfrak{d}\) is isomorphic to \(\overline{D}\) (cf.~Algebra, Chapter VIII, §1, Exercises 9 and 3). Show that, if further B is an Azumaya algebra over A (Chapter II, §5, Exercise 14), so is D.
\end{enumerate}

\begin{enumerate}
\def\labelenumi{\arabic{enumi}.}
\setcounter{enumi}{5}
\tightlist
\item
  Give an example of a non-commutative Artinian ring A, which is filtered by a sequence \((a_{n})\) of two-sided ideals such that \(a_{0} = A\) and \(a_{2} = 0\) and for which Proposition 8 of no. 6 does not hold (cf.~Algebra, Chapter VIII, §2, Exercise 6).
\end{enumerate}

¶ 7. (a) Let A be a commutative ring and m an ideal of A such that A is Hausdorff and complete with the m-adic topology and \(m/m^{2}\) is a finitely generated A-module. Show that the topology on \(A' = A\{X_{1}, \ldots, X_{r}\}\) is the \(m'\) -adic topology, where \(m' = mA'\) and that \(m'/m'^{2}\) is a finitely generated \((A'/m')\) -module (use Proposition 14 of §2, no. 11). In particular, if A is Noetherian, so is \(A'\) .

\begin{enumerate}
\def\labelenumi{(\alph{enumi})}
\setcounter{enumi}{1}
\tightlist
\item
  Let A be a commutative Noetherian ring and m an ideal of A such that A is Hausdorff and complete with the m-adic topology. Let \(u: A \rightarrow B\) be a continuous homomorphism from A to a Hausdorff commutative topological ring B making B into an A-algebra. Show that the following conditions are equivalent:
\end{enumerate}

(α) B is Noetherian, its topology is the mB-adic topology, B is complete and B/mB is a finitely generated algebra over A/m.

(β) B is topologically A-isomorphic to \(\lim B_{n}\) , where \((\mathrm{B}_{n})_{n\geqslant1}\) is an inverse system of discrete A-algebras such that the mappings \(\phi_{nm}:B_{m}\to B_{n}\) for \(m\geqslant n\) are surjective, the kernel of \(\phi_{nm}\) is \(m^{n+1}B_{m}\) and \(B_{1}\) is a finitely generated algebra over A/m.

\((\gamma)\) B is topologically A-isomorphic to a quotient of an algebra of the form \(\mathbf{A}\{\mathbf{X}_1,\dots ,\mathbf{X}_r\}\) by a closed ideal.

(To prove that \((\beta)\) implies \((\gamma)\) , use Proposition 14 of §2, no. 11 and Theorem 1 of §2, no. 8.)

\begin{enumerate}
\def\labelenumi{\arabic{enumi}.}
\setcounter{enumi}{7}
\tightlist
\item
  Let A be a linearly topologized commutative ring. The additive group \(A[[X_{1},\ldots,X_{p}]]\) is identified with the product group \(A^{N^{p}}\) and is given the product topology T.
\end{enumerate}

\begin{enumerate}
\def\labelenumi{(\alph{enumi})}
\item
  Show that \(\mathcal{T}\) is compatible with the ring structure on \(\mathbf{A}[[\mathbf{X}_1,\dots ,\mathbf{X}_p]]\) and is a linear topology for which the mappings \(f\mapsto \partial f / \partial \mathbf{X}_i\) are continuous.
\item
  For every element \(P = \sum_{e} \alpha_{e,P} X^{e}\) of \(A[[X_{1}, \ldots, X_{p}]]\) (notation of §1, no. 6) and every complete Hausdorff topological A-algebra B (§2, Exercise 28), an element \(\mathbf{x} = (x_{1}, \ldots, x_{p}) \in B^{p}\) is said to be substitutable in P if the family \((\alpha_{e,P}\mathbf{x}^{e})\) (where we have written \(x^{e} = x_{1}^{e_{1}} \ldots x_{p}^{e_{p}}\) for \(e = (e_{1}, \ldots, e_{p})\) ) converges to 0 in B with respect to the filter of complements of finite subsets of \(N^{p}\) . This family is then summable and its sum is denoted by \(P(\mathbf{x})\) . Show that the set of formal power series \(P \in A[[X_{1}, \ldots, X_{p}]]\) such that x is substitutable in P is a subring \(S_{x}\) of \(A[[X_{1}, \ldots, X_{p}]]\) and the mapping \(P \mapsto P(\mathbf{x})\) is a homomorphism from \(S_{x}\) to B.
\item
  Show that, if \(\mathbf{x}\) is substitutable in \(\mathbf{P}\) and \(\mathbf{y} = (y_1, \ldots, y_p)\) is an element of \(\mathbf{B}^p\) such that the \(y_i\) are topologically nilpotent for \(1 \leqslant i \leqslant p\) , then \(\mathbf{x} + \mathbf{y}\) is substitutable in \(\mathbf{P}\) .
\end{enumerate}

In particular, if there exists in B a neighbourhood of 0 consisting of topologically nilpotent elements, then, for all \(P \in A[[X_{1}, \ldots, X_{p}]]\) , the set D of \(x \in B^{p}\) which are substitutable in P is open and the mapping \(\mathbf{x} \mapsto \mathrm{P}(\mathbf{x})\) from D to B is continuous.

\begin{enumerate}
\def\labelenumi{(\alph{enumi})}
\setcounter{enumi}{3}
\tightlist
\item
  Suppose from now on that A is Hausdorff and complete and let \(A[[X_{1},\ldots,X_{p}]]\) be given the topology T. For a system of q formal power series \(P_{1},\ldots,P_{q}\) of \(A[[X_{1},\ldots,X_{p}]]\) to be substitutable in a formal power series
\end{enumerate}

\[
\mathrm{Q} \in \mathrm{A} [ [ \mathrm{X} _ {1}, \dots , \mathrm{X} _ {q} ] ],
\]

it is necessary and sufficient that the system \((\mathbf{P}_{1}(0),\ldots,\mathbf{P}_{q}(0))\) be substitutable in Q.

\begin{enumerate}
\def\labelenumi{(\alph{enumi})}
\setcounter{enumi}{4}
\item
  Suppose that \(\mathbf{x}\) is substitutable in each of the \(\mathrm{P}_k\) ( \(1 \leqslant k \leqslant q\) ) and that the system \((\mathrm{P}_1, \ldots, \mathrm{P}_q)\) is substitutable in \(\mathbb{Q}\) . Show that \(\mathbf{x}\) is substitutable in \(\mathbb{Q}(\mathrm{P}_1, \ldots, \mathrm{P}_q)\) , that the system \((\mathrm{P}_1(\mathbf{x}), \ldots, \mathrm{P}_q(\mathbf{x}))\) is substitutable in \(\mathbb{Q}\) and that \(\mathbb{Q}(\mathrm{P}_1(\mathbf{x}), \ldots, \mathrm{P}_q(\mathbf{x})) = ((\mathbb{Q}(\mathrm{P}_1, \ldots, \mathrm{P}_q))(\mathbf{x})\) if further one of the following hypotheses is assumed to be satisfied: ( \(\alpha\) ) there is in \(\mathbb{B}\) an ideal \(n\) such that the ordered pair \((\mathbb{B}, n)\) satisfies Hensel's conditions and the \(x_i\) ( \(1 \leqslant i \leqslant p\) ) belong to \(n\) ; ( \(\beta\) ) the formal power series \(\mathbb{Q}\) is restricted.
\item
  Take \(\mathbf{B} = \mathbf{A} = \mathbf{F}_2\) with the discrete topology, \(\mathbf{P} = \mathbf{X} + \mathbf{X}^2\) , \(\mathbf{Q} = \sum_{n=0}^{\infty} \mathbf{X}^{2^n}\) ; then \(\mathbf{P}\) is substitutable in \(\mathbf{Q}\) , \(x = 1\) is substitutable in \(\mathbf{P}\) and in \(\mathbf{Q}(\mathbf{P})\) and \(\mathbf{P}(x)\) is substitutable in \(\mathbf{Q}\) , but \(\mathbf{Q}(\mathbf{P}(x)) = 0\) and \((\mathbf{Q}(\mathbf{P}))(x) = 1\) .
\end{enumerate}
% semantic-fragments: legacy-input-seam
\relax{}
